\documentclass[11pt, a4paper]{ctexart}

\usepackage{assignment_style}

\begin{document}

% =========================================
% P58. 1
% =========================================
\begin{problem}{P58. 1 分别用 Gauss 消去法、列主元消去法和全主元消去法解下列方程组，结果保留 3 位小数。}
    \textbf{(1)}
    \[
        \begin{cases}
            2x_1 + 6x_2 - 4x_3 = 4, \\
            x_1 + 4x_2 - 5x_3 = 3, \\
            6x_1 - x_2 + 18x_3 = 2;
        \end{cases}
    \]

    \textbf{(2)}
    \[
        \begin{cases}
            2x_1 + x_2 + 2x_3 = 6, \\
            4x_1 + 3x_2 + x_3 = 11, \\
            6x_1 + x_2 + 5x_3 = 13.
        \end{cases}
    \]
\end{problem}

\begin{solution}
    \textbf{(1)}

    \textbf{Gauss 消去法：}
    \[
        \left[
        \begin{array}{ccc|c}
            2 & 6 & -4 & 4 \\
            1 & 4 & -5 & 3 \\
            6 & -1 & 18 & 2
        \end{array}
        \right]
        \to
        \left[
        \begin{array}{ccc|c}
            1 & 3 & -2 & 2 \\
            0 & 1 & -3 & 1 \\
            0 & -19 & 30 & -10
        \end{array}
        \right]
        \to
        \left[
        \begin{array}{ccc|c}
            1 & 3 & -2 & 2 \\
            0 & 1 & -3 & 1 \\
            0 & 0 & -27 & 9
        \end{array}
        \right].
    \]
    回代得
    \[
        x_3 = -\frac{1}{3} = -0.333,\qquad
        x_2 = 1 + 3x_3 = 0,\qquad
        x_1 = 2 + 2x_3 - 3x_2 = \frac{4}{3} = 1.333.
    \]

    \textbf{列主元消去法：}
    第一步交换第 1、3 行：
    \[
        \left[
        \begin{array}{ccc|c}
            2 & 6 & -4 & 4 \\
            1 & 4 & -5 & 3 \\
            6 & -1 & 18 & 2
        \end{array}
        \right]
        \to
        \left[
        \begin{array}{ccc|c}
            6 & -1 & 18 & 2 \\
            1 & 4 & -5 & 3 \\
            2 & 6 & -4 & 4
        \end{array}
        \right].
    \]
    消去并在第二列继续选主元：
    \[
        \left[
        \begin{array}{ccc|c}
            6 & -1 & 18 & 2 \\
            0 & 4.167 & -8 & 2.667 \\
            0 & 6.333 & -10 & 3.333
        \end{array}
        \right]
        \to
        \left[
        \begin{array}{ccc|c}
            6 & -1 & 18 & 2 \\
            0 & 6.333 & -10 & 3.333 \\
            0 & 4.167 & -8 & 2.667
        \end{array}
        \right]
        \to
        \left[
        \begin{array}{ccc|c}
            6 & -1 & 18 & 2 \\
            0 & 6.333 & -10 & 3.333 \\
            0 & 0 & -1.420 & 0.474
        \end{array}
        \right].
    \]
    回代得近似解
    \[
        x_3 \approx -0.334,\qquad
        x_2 \approx -0.001,\qquad
        x_1 \approx 1.336.
    \]

    \textbf{全主元消去法：}
    先交换第 1、3 行，再交换第 1、3 列，此时变量次序变为 $(x_3, x_2, x_1)$：
    \[
        \left[
        \begin{array}{ccc|c}
            2 & 6 & -4 & 4 \\
            1 & 4 & -5 & 3 \\
            6 & -1 & 18 & 2
        \end{array}
        \right]
        \to
        \left[
        \begin{array}{ccc|c}
            18 & -1 & 6 & 2 \\
            -5 & 4 & 1 & 3 \\
            -4 & 6 & 2 & 4
        \end{array}
        \right].
    \]
    继续消去并选取第二步主元：
    \[
        \begin{aligned}
            &\left[
            \begin{array}{ccc|c}
                18 & -1 & 6 & 2 \\
                0 & 3.722 & 2.667 & 3.556 \\
                0 & 5.778 & 3.333 & 4.444
            \end{array}
            \right]
            \to
            \left[
            \begin{array}{ccc|c}
                18 & -1 & 6 & 2 \\
                0 & 5.778 & 3.333 & 4.444 \\
                0 & 3.722 & 2.667 & 3.556
            \end{array}
            \right] \\
            &\to
            \left[
            \begin{array}{ccc|c}
                18 & -1 & 6 & 2 \\
                0 & 5.778 & 3.333 & 4.444 \\
                0 & 0 & 0.520 & 0.693
            \end{array}
            \right].
        \end{aligned}
    \]
    回代时按变量次序 $(x_3, x_2, x_1)$ 计算，得到
    \[
        x_1 \approx 1.333,\qquad
        x_2 \approx 0.000,\qquad
        x_3 \approx -0.333.
    \]

    \vspace{1em}
    \textbf{(2)}

    \textbf{Gauss 消去法：}
    \[
        \left[
        \begin{array}{ccc|c}
            2 & 1 & 2 & 6 \\
            4 & 3 & 1 & 11 \\
            6 & 1 & 5 & 13
        \end{array}
        \right]
        \to
        \left[
        \begin{array}{ccc|c}
            2 & 1 & 2 & 6 \\
            0 & 1 & -3 & -1 \\
            0 & -2 & -1 & -5
        \end{array}
        \right]
        \to
        \left[
        \begin{array}{ccc|c}
            2 & 1 & 2 & 6 \\
            0 & 1 & -3 & -1 \\
            0 & 0 & -7 & -7
        \end{array}
        \right].
    \]
    回代得
    \[
        x_3 = 1,\qquad x_2 = 2,\qquad x_1 = 1.
    \]

    \textbf{列主元消去法：}
    交换第 1、3 行：
    \[
        \left[
        \begin{array}{ccc|c}
            2 & 1 & 2 & 6 \\
            4 & 3 & 1 & 11 \\
            6 & 1 & 5 & 13
        \end{array}
        \right]
        \to
        \left[
        \begin{array}{ccc|c}
            6 & 1 & 5 & 13 \\
            4 & 3 & 1 & 11 \\
            2 & 1 & 2 & 6
        \end{array}
        \right].
    \]
    消去后得
    \[
        \left[
        \begin{array}{ccc|c}
            6 & 1 & 5 & 13 \\
            0 & 3.5 & -2.333 & 2.333 \\
            0 & 0.667 & 0.333 & 1.667
        \end{array}
        \right]
        \to
        \left[
        \begin{array}{ccc|c}
            6 & 1 & 5 & 13 \\
            0 & 3.5 & -2.333 & 2.333 \\
            0 & 0 & 0.778 & 0.778
        \end{array}
        \right].
    \]
    回代仍得
    \[
        x_3 = 1.000,\qquad x_2 = 2.000,\qquad x_1 = 1.000.
    \]

    \textbf{全主元消去法：}
    该方程组在消去过程中无需额外换列，全主元法与列主元法得到同样的上三角系统：
    \[
        \left[
        \begin{array}{ccc|c}
            6 & 1 & 5 & 13 \\
            0 & 3.5 & -2.333 & 2.333 \\
            0 & 0 & 0.778 & 0.778
        \end{array}
        \right],
    \]
    因而
    \[
        x_1 = 1.000,\qquad x_2 = 2.000,\qquad x_3 = 1.000.
    \]
\end{solution}


% =========================================
% 2
% =========================================
\begin{problem}{2. 用三角分解法解题 1 中的两个方程组。}
\end{problem}

\begin{solution}
    \textbf{(1)} 对于
    \[
        A =
        \begin{pmatrix}
            2 & 6 & -4 \\
            1 & 4 & -5 \\
            6 & -1 & 18
        \end{pmatrix},
        \qquad
        b =
        \begin{pmatrix}
            4 \\
            3 \\
            2
        \end{pmatrix},
    \]
    有
    \[
        A = LU =
        \begin{pmatrix}
            1 & 0 & 0 \\
            \frac{1}{2} & 1 & 0 \\
            3 & -19 & 1
        \end{pmatrix}
        \begin{pmatrix}
            2 & 6 & -4 \\
            0 & 1 & -3 \\
            0 & 0 & -27
        \end{pmatrix}.
    \]

    令 $Ly = b$，则
    \[
        y_1 = 4,\qquad
        \frac{1}{2} y_1 + y_2 = 3,\qquad
        3y_1 - 19y_2 + y_3 = 2,
    \]
    解得
    \[
        y =
        \begin{pmatrix}
            4 \\
            1 \\
            9
        \end{pmatrix}.
    \]
    再由 $Ux = y$ 得
    \[
        \begin{pmatrix}
            2 & 6 & -4 \\
            0 & 1 & -3 \\
            0 & 0 & -27
        \end{pmatrix}
        \begin{pmatrix}
            x_1 \\
            x_2 \\
            x_3
        \end{pmatrix}
        =
        \begin{pmatrix}
            4 \\
            1 \\
            9
        \end{pmatrix},
    \]
    从而
    \[
        x_3 = -\frac{1}{3},\qquad x_2 = 0,\qquad x_1 = \frac{4}{3}.
    \]

    \vspace{1em}
    \textbf{(2)} 对于
    \[
        A =
        \begin{pmatrix}
            2 & 1 & 2 \\
            4 & 3 & 1 \\
            6 & 1 & 5
        \end{pmatrix},
        \qquad
        b =
        \begin{pmatrix}
            6 \\
            11 \\
            13
        \end{pmatrix},
    \]
    有
    \[
        A = LU =
        \begin{pmatrix}
            1 & 0 & 0 \\
            2 & 1 & 0 \\
            3 & -2 & 1
        \end{pmatrix}
        \begin{pmatrix}
            2 & 1 & 2 \\
            0 & 1 & -3 \\
            0 & 0 & -7
        \end{pmatrix}.
    \]

    由 $Ly = b$，
    \[
        y_1 = 6,\qquad
        2y_1 + y_2 = 11,\qquad
        3y_1 - 2y_2 + y_3 = 13,
    \]
    解得
    \[
        y =
        \begin{pmatrix}
            6 \\
            -1 \\
            -7
        \end{pmatrix}.
    \]
    再由 $Ux = y$，
    \[
        \begin{pmatrix}
            2 & 1 & 2 \\
            0 & 1 & -3 \\
            0 & 0 & -7
        \end{pmatrix}
        \begin{pmatrix}
            x_1 \\
            x_2 \\
            x_3
        \end{pmatrix}
        =
        \begin{pmatrix}
            6 \\
            -1 \\
            -7
        \end{pmatrix},
    \]
    得
    \[
        x_3 = 1,\qquad x_2 = 2,\qquad x_1 = 1.
    \]
\end{solution}


% =========================================
% 3
% =========================================
\begin{problem}{3. 设线性方程组}
    \[
        \begin{pmatrix}
            1 & 2 & 3 & 4 \\
            1 & 4 & 9 & 16 \\
            1 & 8 & 27 & 64 \\
            1 & 16 & 81 & 256
        \end{pmatrix}
        x
        =
        \begin{pmatrix}
            2 \\
            10 \\
            44 \\
            190
        \end{pmatrix}.
    \]
    用三角分解法求解。
\end{problem}

\begin{solution}
    对系数矩阵作 $LU$ 分解可得
    \[
        L =
        \begin{pmatrix}
            1 & 0 & 0 & 0 \\
            1 & 1 & 0 & 0 \\
            1 & 3 & 1 & 0 \\
            1 & 7 & 6 & 1
        \end{pmatrix},
        \qquad
        U =
        \begin{pmatrix}
            1 & 2 & 3 & 4 \\
            0 & 2 & 6 & 12 \\
            0 & 0 & 6 & 24 \\
            0 & 0 & 0 & 24
        \end{pmatrix},
    \]
    从而 $A = LU$。

    先解下三角方程 $Ly = b$：
    \[
        \begin{cases}
            y_1 = 2, \\
            y_1 + y_2 = 10, \\
            y_1 + 3y_2 + y_3 = 44, \\
            y_1 + 7y_2 + 6y_3 + y_4 = 190,
        \end{cases}
    \]
    得
    \[
        y =
        \begin{pmatrix}
            2 \\
            8 \\
            18 \\
            24
        \end{pmatrix}.
    \]

    再解上三角方程 $Ux = y$：
    \[
        \begin{cases}
            24x_4 = 24, \\
            6x_3 + 24x_4 = 18, \\
            2x_2 + 6x_3 + 12x_4 = 8, \\
            x_1 + 2x_2 + 3x_3 + 4x_4 = 2,
        \end{cases}
    \]
    得
    \[
        x =
        \begin{pmatrix}
            -1 \\
            1 \\
            -1 \\
            1
        \end{pmatrix}.
    \]
\end{solution}


% =========================================
% 6
% =========================================
\begin{problem}{6. 分别用平方根法和改进平方根法求解方程组}
    \[
        \begin{pmatrix}
            1 & 2 & 1 & -3 \\
            2 & 5 & 0 & -5 \\
            1 & 0 & 14 & 1 \\
            -3 & -5 & 1 & 15
        \end{pmatrix}
        x
        =
        \begin{pmatrix}
            1 \\
            2 \\
            16 \\
            8
        \end{pmatrix}.
    \]
\end{problem}

\begin{solution}
    由于系数矩阵 $A$ 对称正定，故可分别进行 $LL^{\mathrm{T}}$ 分解和 $LDL^{\mathrm{T}}$ 分解。

    \textbf{(1) 平方根法（Cholesky 分解）}

    先取
    \[
        l_{11} = \sqrt{1} = 1,\qquad
        l_{11} L_{21} =
        \begin{pmatrix}
            2 \\
            1 \\
            -3
        \end{pmatrix},
    \]
    从而
    \[
        L_{21} =
        \begin{pmatrix}
            2 \\
            1 \\
            -3
        \end{pmatrix}.
    \]
    余下主子块满足
    \[
        \begin{pmatrix}
            5 & 0 & -5 \\
            0 & 14 & 1 \\
            -5 & 1 & 15
        \end{pmatrix}
        -
        \begin{pmatrix}
            4 & 2 & -6 \\
            2 & 1 & -3 \\
            -6 & -3 & 9
        \end{pmatrix}
        =
        \begin{pmatrix}
            1 & -2 & 1 \\
            -2 & 13 & 4 \\
            1 & 4 & 6
        \end{pmatrix}.
    \]
    继续分解该 $3 \times 3$ 块，可得
    \[
        L =
        \begin{pmatrix}
            1 & 0 & 0 & 0 \\
            2 & 1 & 0 & 0 \\
            1 & -2 & 3 & 0 \\
            -3 & 1 & 2 & 1
        \end{pmatrix},
        \qquad
        A = LL^{\mathrm{T}}.
    \]

    先解 $Ly = b$：
    \[
        y =
        \begin{pmatrix}
            1 \\
            0 \\
            5 \\
            1
        \end{pmatrix}.
    \]
    再解 $L^{\mathrm{T}}x = y$，得
    \[
        x =
        \begin{pmatrix}
            1 \\
            1 \\
            1 \\
            1
        \end{pmatrix}.
    \]

    \vspace{1em}
    \textbf{(2) 改进平方根法（$LDL^{\mathrm{T}}$ 分解）}

    取单位下三角矩阵
    \[
        L =
        \begin{pmatrix}
            1 & 0 & 0 & 0 \\
            2 & 1 & 0 & 0 \\
            1 & -2 & 1 & 0 \\
            -3 & 1 & \frac{2}{3} & 1
        \end{pmatrix},
        \qquad
        D = \operatorname{diag}(1, 1, 9, 1),
    \]
    则
    \[
        A = LDL^{\mathrm{T}}.
    \]

    先解 $Lz = b$：
    \[
        z =
        \begin{pmatrix}
            1 \\
            0 \\
            15 \\
            1
        \end{pmatrix}.
    \]
    再解 $Dw = z$：
    \[
        w =
        \begin{pmatrix}
            1 \\
            0 \\
            \frac{5}{3} \\
            1
        \end{pmatrix}.
    \]
    最后解 $L^{\mathrm{T}}x = w$，得
    \[
        x =
        \begin{pmatrix}
            1 \\
            1 \\
            1 \\
            1
        \end{pmatrix}.
    \]
\end{solution}


% =========================================
% 8
% =========================================
\begin{problem}{8. 用追赶法解下列方程组}
    \[
        \begin{pmatrix}
            2 & -1 & 0 & 0 \\
            -1 & 2 & -1 & 0 \\
            0 & -1 & 2 & -1 \\
            0 & 0 & -1 & 2
        \end{pmatrix}
        x
        =
        \begin{pmatrix}
            0 \\
            1 \\
            0 \\
            2.5
        \end{pmatrix}.
    \]
\end{problem}

\begin{solution}
    系数矩阵为三对角矩阵，可作如下分解：
    \[
        A = LU,
        \qquad
        L =
        \begin{pmatrix}
            1 & 0 & 0 & 0 \\
            -\frac{1}{2} & 1 & 0 & 0 \\
            0 & -\frac{2}{3} & 1 & 0 \\
            0 & 0 & -\frac{3}{4} & 1
        \end{pmatrix},
        \qquad
        U =
        \begin{pmatrix}
            2 & -1 & 0 & 0 \\
            0 & \frac{3}{2} & -1 & 0 \\
            0 & 0 & \frac{4}{3} & -1 \\
            0 & 0 & 0 & \frac{5}{4}
        \end{pmatrix}.
    \]

    先解 $Ly = b$：
    \[
        \begin{cases}
            y_1 = 0, \\
            -\dfrac{1}{2} y_1 + y_2 = 1, \\
            -\dfrac{2}{3} y_2 + y_3 = 0, \\
            -\dfrac{3}{4} y_3 + y_4 = 2.5,
        \end{cases}
    \]
    得
    \[
        y =
        \begin{pmatrix}
            0 \\
            1 \\
            \frac{2}{3} \\
            3
        \end{pmatrix}.
    \]

    再解 $Ux = y$：
    \[
        \begin{cases}
            \dfrac{5}{4}x_4 = 3, \\
            \dfrac{4}{3}x_3 - x_4 = \dfrac{2}{3}, \\
            \dfrac{3}{2}x_2 - x_3 = 1, \\
            2x_1 - x_2 = 0,
        \end{cases}
    \]
    得
    \[
        x_4 = \frac{12}{5},\qquad
        x_3 = \frac{23}{10},\qquad
        x_2 = \frac{11}{5},\qquad
        x_1 = \frac{11}{10}.
    \]
    即
    \[
        x =
        \begin{pmatrix}
            \frac{11}{10} \\
            \frac{11}{5} \\
            \frac{23}{10} \\
            \frac{12}{5}
        \end{pmatrix}.
    \]
\end{solution}


% =========================================
% 12
% =========================================
\begin{problem}{12. 已知病态方程组}
    \[
        \begin{cases}
            x_1 + 0.99x_2 = 1, \\
            0.99x_1 + 0.98x_2 = 1,
        \end{cases}
        \qquad
        \text{其精确解为 }
        x =
        \begin{pmatrix}
            100 \\
            -100
        \end{pmatrix}.
    \]
    \begin{enumerate}[1.]
        \item 计算系数矩阵的条件数 $\operatorname{cond}_{\infty}(A)$；
        \item 取 $x_1^* = (1, 0)^{\mathrm{T}}$，$x_2^* = (100.5, -99.5)^{\mathrm{T}}$，分别计算残量 $r_i = b - Ax_i^*$，并说明结果说明了什么。
    \end{enumerate}
\end{problem}

\begin{solution}
    记
    \[
        A =
        \begin{pmatrix}
            1 & 0.99 \\
            0.99 & 0.98
        \end{pmatrix},
        \qquad
        b =
        \begin{pmatrix}
            1 \\
            1
        \end{pmatrix}.
    \]

    \textbf{(1) 计算条件数}

    \[
        A^{-1}
        =
        \frac{1}{1 \cdot 0.98 - 0.99 \cdot 0.99}
        \begin{pmatrix}
            0.98 & -0.99 \\
            -0.99 & 1
        \end{pmatrix}
        =
        \begin{pmatrix}
            -9800 & 9900 \\
            9900 & -10000
        \end{pmatrix}.
    \]
    因而
    \[
        \|A\|_{\infty} = \max\{1 + 0.99,\ 0.99 + 0.98\} = 1.99,
    \]
    \[
        \|A^{-1}\|_{\infty} = \max\{9800 + 9900,\ 9900 + 10000\} = 19900.
    \]
    所以
    \[
        \operatorname{cond}_{\infty}(A)
        = \|A\|_{\infty}\|A^{-1}\|_{\infty}
        = 1.99 \times 19900
        = 39601.
    \]

    \textbf{(2) 比较残量与误差}

    对于近似解
    \[
        x_1^* =
        \begin{pmatrix}
            1 \\
            0
        \end{pmatrix},
    \]
    有
    \[
        r_1 = b - Ax_1^*
        =
        \begin{pmatrix}
            1 \\
            1
        \end{pmatrix}
        -
        \begin{pmatrix}
            1 \\
            0.99
        \end{pmatrix}
        =
        \begin{pmatrix}
            0 \\
            0.01
        \end{pmatrix}.
    \]
    其相对误差为
    \[
        \frac{\|x - x_1^*\|_{\infty}}{\|x\|_{\infty}}
        =
        \frac{100}{100}
        = 1.
    \]

    对于近似解
    \[
        x_2^* =
        \begin{pmatrix}
            100.5 \\
            -99.5
        \end{pmatrix},
    \]
    有
    \[
        Ax_2^*
        =
        \begin{pmatrix}
            1.995 \\
            1.985
        \end{pmatrix},
        \qquad
        r_2 = b - Ax_2^*
        =
        \begin{pmatrix}
            -0.995 \\
            -0.985
        \end{pmatrix}.
    \]
    其相对误差为
    \[
        \frac{\|x - x_2^*\|_{\infty}}{\|x\|_{\infty}}
        =
        \frac{0.5}{100}
        = 5 \times 10^{-3}.
    \]

    虽然 $\|r_1\|_{\infty} < \|r_2\|_{\infty}$，但 $x_1^*$ 的相对误差却远大于 $x_2^*$。这说明对于病态方程组，仅凭残量大小并不能可靠反映近似解的误差，因为此时 $\operatorname{cond}(A) \gg 1$，微小扰动也会被显著放大。
\end{solution}


% =========================================
% 13
% =========================================
\begin{problem}{13. 求超定方程组的最小二乘解}
    \[
        \begin{cases}
            2x_1 + 4x_2 = 11, \\
            3x_1 - 5x_2 = 3, \\
            x_1 + 2x_2 = 6, \\
            2x_1 + x_2 = 7.
        \end{cases}
    \]
\end{problem}

\begin{solution}
    记
    \[
        A =
        \begin{pmatrix}
            2 & 4 \\
            3 & -5 \\
            1 & 2 \\
            2 & 1
        \end{pmatrix},
        \qquad
        b =
        \begin{pmatrix}
            11 \\
            3 \\
            6 \\
            7
        \end{pmatrix}.
    \]
    最小二乘解满足正规方程
    \[
        A^{\mathrm{T}} A x = A^{\mathrm{T}} b.
    \]
    计算得
    \[
        A^{\mathrm{T}}A =
        \begin{pmatrix}
            18 & -3 \\
            -3 & 46
        \end{pmatrix},
        \qquad
        A^{\mathrm{T}}b =
        \begin{pmatrix}
            51 \\
            48
        \end{pmatrix}.
    \]
    又因为
    \[
        |A^{\mathrm{T}}A| = 18 \cdot 46 - (-3)(-3) = 819 \ne 0,
    \]
    所以正规方程有唯一解。进一步有
    \[
        (A^{\mathrm{T}}A)^{-1}
        =
        \frac{1}{819}
        \begin{pmatrix}
            46 & 3 \\
            3 & 18
        \end{pmatrix}.
    \]
    因而
    \[
        x
        =
        (A^{\mathrm{T}}A)^{-1}A^{\mathrm{T}}b
        =
        \frac{1}{819}
        \begin{pmatrix}
            46 & 3 \\
            3 & 18
        \end{pmatrix}
        \begin{pmatrix}
            51 \\
            48
        \end{pmatrix}
        =
        \frac{1}{819}
        \begin{pmatrix}
            2490 \\
            1017
        \end{pmatrix}
        =
        \begin{pmatrix}
            \dfrac{830}{273} \\
            \dfrac{113}{91}
        \end{pmatrix}.
    \]
\end{solution}

\end{document}
